\scanpage{65}
\mspage{123}
\begin{flushright}n° 371\end{flushright}

$(i_h)_{1\handwritten{\leq}h\handwritten{\leq}p}$ est la suite strictement croissante des éléments de $H$.

PROPOSITION 6.- \uline{Soit} $\unsure{K|M}$ (resp. $N$) \uline{un} $A$-\uline{module ayant une base} $(e_i)_{1\handwritten{\leq}i\handwritten{\leq}n}$ (resp. $(f_j)_{1\handwritten{\leq}j\handwritten{\leq}m}$) \uline{de} $n$ (resp. $m$) \uline{éléments} ; \uline{soit} (pour $p\handwritten{\leq}\inf(m,n)$) $(e_K)$ (resp. $(f_H)$) \uline{la base correspondante de} $\overset{p}{\handwritten{\bigwedge}}M$ (resp. $\overset{p}{\handwritten{\bigwedge}}N$). \uline{Soit} $u:M\handwritten{\to}N$ \uline{une application} $A$-\uline{linéaire}, \uline{et désignons par} $\underline{X}$ \uline{la matrice de} $u$ \uline{par rapport aux bases} $(e_i)$ \uline{et} $(f_j)$. \uline{Alors la matrice de} $\overset{p}{\handwritten{\bigwedge}}u$ \uline{par rapport aux bases} $(e_K)$ \uline{et} $(f_H)$ \uline{est la matrice} $(X_{H,K})$ \uline{à} $\binom{m}{p}$ \uline{lignes et} $\binom{n}{p}$ \uline{colonnes formée des mineurs d'ordre} $p$ \uline{de} $\underline{X}$.

En effet soit $(j_k)_{1\handwritten{\leq}k\handwritten{\leq}p}$ la suite strictement croissante des éléments de $K\handwritten{\subset}\handwritten{[}1,n\handwritten{]}$ ; on a par définition
\[
\overset{p}{\handwritten{\bigwedge}}u(e_K)=u(e_{j_1})\handwritten{\wedge}u(e_{j_2})\handwritten{\wedge}\ldots\handwritten{\wedge}u(e_{j_p})\ ;
\tag*{(10)}
\]
donc l'élément de la matrice de $\overset{p}{\handwritten{\bigwedge}}u$ qui se trouve dans la ligne d'indice $H$ et la colonne d'indice $K$ est la composante d'indice $H$ du $p$-vecteur décomposable $u(e_{j_1})\handwritten{\wedge}\ldots\handwritten{\wedge}u(e_{j_p})$ ; il est par suite égal à $X_{H,K}$ en vertu de la prop.5.

On dit que la matrice $(X_{H,K})$ est la \uline{puissance extérieure} $p$-\uline{ème} de $\underline{X}$ et on la note $\overset{p}{\handwritten{\bigwedge}}\underline{X}$.

4. \uline{Développements d'un déterminant}.

Reprenons la formule (4) du n°1 donnant le déterminant d'une matrice carrée $\underline{X}=(\handwritten{\xi}_{ij})$ d'ordre $n$ ; pour toute partie $H$ de l'ensemble $\handwritten{[}1,n\handwritten{]}$, désignons par $H'$ la partie complémentaire de $H$ dans $\handwritten{[}1,n\handwritten{]}$ ; si $(i_k)_{1\handwritten{\leq}i\handwritten{\leq}p}$ (resp. $(j_h)_{1\handwritten{\leq}h\handwritten{\leq}n-p}$) est la suite strictement croissante des éléments de $H$ (resp. $H'$), on peut écrire (\S\ 3, n°2, lemme 3)
\[
x_1\handwritten{\wedge}x_2\handwritten{\wedge}\ldots\handwritten{\wedge}x_n
=\handwritten{\rho}_{H,H'}(x_{i_1}\handwritten{\wedge}x_{i_2}\handwritten{\wedge}\ldots\handwritten{\wedge}x_{i_p})\handwritten{\wedge}(x_{j_1}\handwritten{\wedge}x_{j_2}\handwritten{\wedge}\ldots\handwritten{\wedge}x_{j_{n-p}})
\tag*{(11)}
\]
où $\handwritten{\rho}_{H,H'}$ a la signification définie au \S\ 5, n°3, formule (12). Avec les notations du n°3, on a
% archive sheet number: 64
